%
% File: appendices.tex
% Author: Diar Begisbayev
%
\chapter{Appendix}
\label{app:appendix}
\setlength{\parskip}{1em}

\section{Public Holidays and Special Events}
\label{app:holidays}

Table~\ref{Tab.holidays} lists the public holidays and special events included in the synthetic dataset for the year 2025. These events were modelled with attendance-driven demand multipliers ranging from 5{,}000 to 60{,}000 attendees.

\begin{table}[H]
\centering
\caption{Public holidays and special events in the synthetic dataset (2025).}
\label{Tab.holidays}
\begin{tabular}{llcc}
\toprule
Date & Event & Type & Attendance \\
\midrule
1 January & New Year & Public Holiday & --- \\
7 January & Orthodox Christmas & Public Holiday & --- \\
8 March & International Women's Day & Public Holiday & --- \\
21--23 March & Nauryz Meyrami & Public Holiday & --- \\
1 May & Kazakhstan Unity Day & Public Holiday & --- \\
7 May & Defender of the Fatherland Day & Public Holiday & --- \\
9 May & Victory Day & Public Holiday & --- \\
6 July & Capital City Day & Public Holiday & --- \\
30 August & Constitution Day & Public Holiday & --- \\
25 October & Republic Day & Public Holiday & --- \\
16 December & Independence Day & Public Holiday & --- \\
15--20 June & EXPO-2025 Astana & Special Event & 40{,}000 \\
10 August & City Marathon & Special Event & 12{,}000 \\
5 September & Astana Day Concert & Special Event & 25{,}000 \\
\bottomrule
\end{tabular}
\end{table}

\section{Network Route Definitions}
\label{app:routes}

Table~\ref{Tab.routes} summarises the ten synthetic bus routes generated from OpenStreetMap data for the Astana network.

\begin{table}[H]
\centering
\caption{Synthetic bus route definitions.}
\label{Tab.routes}
\begin{tabular}{llcc}
\toprule
Route ID & Orientation & Stops & Districts Served \\
\midrule
R1 & North--South & 18 & Esil, Saryarka \\
R2 & East--West & 20 & Almaty, Baikonur \\
R3 & Central Loop & 16 & Esil \\
R4 & Peripheral Ring & 14 & All four \\
R5 & Express Corridor & 12 & Esil, Almaty \\
R6 & North Loop & 15 & Saryarka \\
R7 & South Cross & 17 & Baikonur, Almaty \\
R8 & University Line & 19 & Esil, Saryarka \\
R9 & Residential Line & 20 & Baikonur \\
R10 & Business District & 14 & Esil \\
\bottomrule
\end{tabular}
\end{table}

\section{Feature Engineering Details}
\label{app:features}

Table~\ref{Tab.features} provides the complete list of 16 input features with their descriptions, ranges, and normalization strategies.

\begin{table}[H]
\centering
\caption{Complete input feature specifications.}
\label{Tab.features}
\begin{tabular}{cl>{\raggedright\arraybackslash}p{4.5cm}c@{}}
\toprule
Idx & Feature & Description & Normalisation \\
\midrule
0 & passengers\_boarding & Hourly passenger boardings at station & z-score \\
1 & passengers\_alighting & Hourly passenger alightings at station & z-score \\
2 & load & Net passengers on vehicle after stop & z-score \\
3 & temperature & Ambient temperature ($^\circ$C) & z-score \\
4 & precipitation & Binary rain/snow indicator & none \\
5 & is\_holiday & Binary public holiday flag & none \\
6--7 & hour\_sin, hour\_cos & Diurnal cyclical encoding & z-score \\
8--9 & dow\_sin, dow\_cos & Day-of-week cyclical encoding & z-score \\
10 & rush\_hour & Binary morning/evening peak flag & none \\
11 & delta\_h & Hours to next holiday & z-score \\
12 & roll\_6h & 6-hour rolling mean of boardings & z-score \\
13 & roll\_24h & 24-hour rolling mean of boardings & z-score \\
14 & dev\_24h & Deviation from 24-hour rolling mean & z-score \\
15 & ratio\_24h & Ratio of current to 24-hour mean & z-score \\
\bottomrule
\end{tabular}
\end{table}

\section{Hyperparameter Search Space}
\label{app:hyperparams}

Table~\ref{Tab.hypersearch} reports the hyperparameter search space explored during model development. Final values are shown in bold.

\begin{table}[H]
\centering
\caption{Hyperparameter search space and final values.}
\label{Tab.hypersearch}
\begin{tabular}{lccc}
\toprule
Parameter & Search Range & Final Value \\
\midrule
Model dimension $d_{\text{model}}$ & \{128, 192, 256\} & \textbf{192} \\
LoRA rank $r$ & \{8, 16, 32\} & \textbf{16} \\
Attention heads $n_h$ & \{4, 6, 8\} & \textbf{6} \\
Graph hops $K$ & \{2, 3, 4\} & \textbf{3} \\
Learning rate & \{1e-4, 3e-4, 1e-3\} & \textbf{3e-4} \\
Weight decay & \{1e-4, 1e-3, 1e-2\} & \textbf{1e-3} \\
Dropout & \{0.0, 0.1, 0.2\} & \textbf{0.1} \\
Auxiliary loss $\lambda$ & \{0.1, 0.3, 0.5\} & \textbf{0.3} \\
Batch size & \{16, 32, 64\} & \textbf{32} \\
Warmup epochs & \{10, 20, 30\} & \textbf{20} \\
\bottomrule
\end{tabular}
\end{table}

\section{Software and Hardware Environment}
\label{app:environment}

All experiments were conducted on a single workstation with the following specifications:

\begin{itemize}
    \item \textbf{CPU}: AMD Ryzen 5 5600X (6 cores, 12 threads)
    \item \textbf{GPU}: NVIDIA GeForce RTX 3060 (12 GB GDDR6)
    \item \textbf{RAM}: 32 GB DDR4-3200
    \item \textbf{OS}: Windows 11 Pro (build 26200)
    \item \textbf{Python}: 3.11.8
    \item \textbf{PyTorch}: 2.12.0+cu126
    \item \textbf{CUDA}: 12.6
    \item \textbf{Key packages}: NumPy 1.26, pandas 2.1, scikit-learn 1.4, matplotlib 3.8, networkx 3.2
\end{itemize}

The complete source code, training scripts, and model weights will be released upon acceptance to facilitate reproducibility.
